\section{Campos con radios dados por las fibras}
Fijemos $v=|U|\le L\log n$ y el exterior común sin esas $v$ marcas. Su núcleo balanceado
$z$ y sus márgenes tienen cotas uniformes por \emph{Grados intermedios}, parte 04:
quitar las marcas cuesta $O(v/j)$ por margen. Para $i\in U$ escribimos
\[
 d_i=\max_{f\ni i}|f\cap U|,\qquad
 \|x\|_*=\max_i d_i |x_i|,\qquad \sum_{i\in f\cap U}1/d_i\le1 .
\]
Aquí $f$ recorre las cuatro familias, contando también las fibras unitarias.
Se desdoblan a lo sumo $2v$ índices RC, con masa $s$ para cada copia y masa basal
menos $s$ en el original. El parámetro $s$ vale $1/j$ en el tablero. En cada copia
sea $t_b=1-\sum_{i\in b}x_i$, y en índices ordinarios $t_b=1$. Si $\|x\|_*\le a$,
$|t_b-1|\le a$. La masa $M_1=t\mu$ difiere en TV del sistema basal, identificando
copias, en $O(v|s|)$. Ambas masas totales son iguales. El equilibrio exterior es
\[
 w(1-st_rt_c w)=z e^{\alpha_r+\beta_c},\qquad
 \int w\,dM_1^C=\int w\,dM_1^R=1 .
\]
En calibre fijo el inverso del jacobiano es acotado en supremo por brecha y núcleo
acotado: la ecuación coordenada a coordenada convierte el control en norma dos en
control en supremo. Las copias de masa cero tienen bloque diagonal identidad.
El argumento de función implícita de la parte 04 citada vale para $v|s|\le\kappa$,
con $a,\kappa$ fijos pequeños, pues sólo intervienen TV, $\sup|t|$ y el radio de
la raíz en cero. Las ecuaciones son holomorfas también en ceros estructurales por $w/z$.
\paragraph{Derivada de la acción estacionaria.}
Escribamos $u=t_rt_cw$ y
\[
 E_s(u;z)=-u\log(u/z)-u+\sum_{\ell\ge2}\frac{s^{\ell-1}u^\ell}{\ell(\ell-1)}.
\]
La acción exterior sin su campo basal es $\mathcal R/s$, donde
\[
 \mathcal R=\int E_s(u;z)\,d\mu_R d\mu_C+
        \int t\log t\,d\mu_R+\int t\log t\,d\mu_C .
\]
Esto se obtiene sustituyendo $u$ y aplicando ambos márgenes. En la silla,
$E_s'=-\log t_r-\log t_c-\alpha_r-\beta_c-2$. Además $\int u\,d\mu_C=t_r$
y simétricamente. Al variar $x_i$, dos márgenes varían en $-1$ cada uno con masa $s$.
Al integrar $E'_s\,du$, los dos términos $\log t$ y las dos unidades se cancelan
con la variación de $t\log t$. Resulta, por continuidad también en $s=0$,
\[
 \partial_i(\mathcal R/s)=\alpha_{r_i}+\beta_{c_i}.
\]
Sea $R_c$ esta acción tras sustraer su valor y campo afín en $x=0$ para cada $s$.
Su gradiente es acotado en la caja, y en media caja
$\sum_l|\partial_l\partial_i R_c|/d_l\le C$. En efecto, derívese en una dirección
$h_l=e^{\mathrm i\theta_l}/d_l$ alineando las fases en la fila $i$ y aplíquese
Cauchy en el parámetro escalar. Esto evita una pérdida del número de vecinos.
\paragraph{Silla y determinante locales.}
Con $d_i^{\rm inc}\le4$ factores duros activos por sitio y $q_f=1-\sum_{i\in f}x_i$,
la ecuación local es
\[
 x_i=\zeta_i e^{\partial_i R_c(x)}
         \prod_{f\ni i} q_f/(1-x_i)^{d_i^{\rm inc}-1}.
\]
Para $|\zeta_i|\le\kappa'/d_i$, la derecha y su derivada en norma ponderada
cuestan $C\kappa'$. Se obtiene la rama estacionaria en la caja por contracción.
La Hessiana local negativa se multiplica a derecha por
$\operatorname{diag}(x_i(1-x_i))$. En la matriz resultante $I+\Delta$,
las incidencias de cada fibra cuestan $C\sum_{i\in f}|x_i|\le C\|x\|_*$
y la parte exterior cuesta $C\|x\|_*$ por la cota del gradiente diferenciado.
Así $|\log\det(I+\Delta)|\le Cv\|\Delta\|_\infty$ en una caja reducida.
Los demás logaritmos son los de $q_f$, $1-x_i$ y $x_i/\zeta_i$, este último
prolongado por la ecuación de silla. La acción centrada se integra a partir de su
gradiente. El prefactor exterior es acotado como en la parte siguiente, también aquí
con $|t_b|\le2$. El valor canónico sin constantes comunes es $O(v)$ en el polidisco.
En los campos físicos $\zeta_i=O(1/j)$ aplicando la ecuación con sólo RC y $x_i=P_i$.
